% =====================================================================
% 1. Introduction
% =====================================================================
\section{Introduction}
\label{sec:introduction}

\subsection{Background}

When a low-altitude small UAV enters the surveillance airspace, the reconnaissance system can only provide a coarse alert---an approximate azimuth, horizontal distance, and altitude---which does not give a precise position. The nominal alert values are therefore used only as planning data (Assumption~A5): they fix the target position for feasibility checks, station selection, and scheduling, but are neither a measurement nor a localization result. The precise horizontal position is obtained, as Problem~2 prescribes, by two-station synchronous direction finding: a single vehicle yields only one bearing line, while two vehicles observing synchronously from different directions intersect two bearing lines. The target altitude does not enter the intersection and enters only through elevation constraints, building occlusion, and pointing errors. The accuracy analysis of Problem~2 is thus a conditional analysis of this bearing-intersection process, whose only stochastic input is the specified azimuth error. A baseline that is too short, or an intersection angle that is too small or too large, amplifies the error; station geometry is therefore the key to accuracy.

\subsection{Problems to Be Solved}

\begin{enumerate}[label=(\arabic*)]
\item \textbf{Single-target observable region and optimal station.} For T1, determine and plot the observable region under the vehicle-geometry, LOS-occlusion and gimbal-pitch constraints. Using $Q=1/r$, find the optimal station, distance and elevation, and state the optimality criterion and tie-breaking rules.

\item \textbf{Two-station intersection accuracy and multi-target assignment.} Under azimuth errors $\mathcal{N}(0,(1^\circ)^2)$, relate accuracy to baseline, intersection angle and range, and define an accuracy index. Under a 30~s travel limit, select two of five vehicles for T2 and analyze degeneracy; then minimize the latest observation start for three alerts on six slots shared by five vehicles, output the timeline, and define a nearest-assignment baseline.

\item \textbf{Localization and pointing error propagation.} With $\sigma(s)=0.3+0.02s$~m, gimbal bias $0.8^\circ$, control error $0.6^\circ$ and a 40~m range, build the $\Delta\psi$ propagation model and find the maximum travel distance satisfying $P(|\Delta\psi|\le5^\circ)\ge0.95$. Derive the $(a,b)$ constraints at 80~m, justify the approximation, and compare the three error sources.

\item \textbf{Rescheduling under disturbances.} Taking the 2(2) plan as the baseline, at $t=5$~s consider a new alert ($\theta\in[0^\circ,360^\circ)$, $\rho\in[30,50]$~m, $H=50$~m), one vehicle failure, an outage $[5,t_{\mathrm{rec}}]$ ($t_{\mathrm{rec}}\in(5,30]$) and their combinations. Build a rescheduling model, lock started observations and allow redispatch, design representative scenarios, and report $(\eta,t_{\mathrm{last}})$ within $[0,60]$~s.
\end{enumerate}

% =====================================================================
% 2. Problem Analysis
% =====================================================================
\section{Problem Analysis}
\label{sec:analysis}

The four problems progress along the chain ``geometric feasibility---intersection accuracy---error propagation---disturbance response'': Q1 answers where observation is possible and where it is best; Q2 extends single-station geometry to two-station intersection and multi-alert scheduling; Q3 advances deterministic geometry to error probability; Q4 advances static scheduling to disturbance response. The feasible region, accuracy conclusions, and timeline of earlier problems form the input baselines of later ones. Figure~\ref{fig:framework} shows the input--output relations and the modeling chain.

\begin{figure}[htbp]
  \centering
  \includegraphics[width=\figwidth]{fig0_framework.png}
  \caption{Framework of the four problems and the overall modeling chain}
  \label{fig:framework}
\end{figure}

\subsection{Geometry and Parameters}

The site is a $100~\mathrm{m}\times100~\mathrm{m}$ plane with the southwest corner as the origin, the $x$-axis east and the $y$-axis north; the dispatch center is $O=(50,50)$. The site contains five axis-aligned cuboid buildings and five standby observation vehicles; their geometry and initial positions are listed in Table~\ref{tab:2-1}.

\begin{table}[H]
\centering
\caption{Building geometry and initial positions of the observation vehicles}
\label{tab:2-1}
{\small
\renewcommand{\arraystretch}{1.08}
\begin{tabularnb}{@{}p{20mm}p{12mm}p{40mm}p{42mm}p{18mm}@{}}
  \toprule
  Category & ID & Planar range (m) & Initial position $(x,y)$ (m) & Height (m) \\
  \midrule
  Building & B1 & $[22,34]\times[16,26]$ & --- & 20 \\
  Building & B2 & $[63,79]\times[59,71]$ & --- & 28 \\
  Building & B3 & $[10,24]\times[65,76]$ & --- & 12 \\
  Building & B4 & $[72,86]\times[22,34]$ & --- & 24 \\
  Building & B5 & $[41,51]\times[38,46]$ & --- & 32 \\
  Vehicle & V1 & --- & $(15,14)$ & --- \\
  Vehicle & V2 & --- & $(88,10)$ & --- \\
  Vehicle & V3 & --- & $(14,86)$ & --- \\
  Vehicle & V4 & --- & $(88,68)$ & --- \\
  Vehicle & V5 & --- & $(50,93)$ & --- \\
  \bottomrule
\end{tabularnb}
}
\par\smallskip
\begin{minipage}{\linewidth}
\noindent\footnotesize Note: All buildings are axis-aligned cuboids standing on $z=0$; a vehicle is an equivalent disk of radius $1.5$~m, and the tabulated position is the initial coordinate of its center (the ground projection of the camera). The three baseline alerts are given in Section~2.1; the complete parameter table is provided in \texttt{params\_table.md} in the supporting materials.
\end{minipage}
\end{table}

A vehicle is modeled as a disk of radius $1.5$~m, travels at a constant $3$~m/s, carries a camera $2$~m above the ground, and has a gimbal pitch range of $[-10^\circ,80^\circ]$ with horizontal and pitch angular rates of $30^\circ$/s; the slew time is the longer axis. During travel the gimbal remains level facing east, and it may turn toward the target only after arrival. A task consists of travel, slew, synchronized $5$~s observation by two vehicles, and reset; all phases are counted on the timeline. A line of sight tangent to a building is judged occluded. An alert is given as $(\theta,R,H)$, and the target's horizontal position is $(50+R\cos\theta,50+R\sin\theta)$; the baseline alerts are T1$(30^\circ,40~\mathrm{m},45~\mathrm{m})$, T2$(150^\circ,45~\mathrm{m},35~\mathrm{m})$, and T3$(285^\circ,38~\mathrm{m},55~\mathrm{m})$ (see~\cite{swmcm}); Problem~4 adds T4$(\theta\in[0^\circ,360^\circ),R\in[30,50]~\mathrm{m},H=50~\mathrm{m})$. These nominal alert values are planning data for all geometry, feasibility, and scheduling computations (Assumption~A5); every accuracy conclusion is conditional on the bearing and pointing error models of Problems~2 and 3.

\subsection{Analysis of Problem 1}

Problem~1 asks where observation is possible, where it is best, how ``optimal'' is defined, and how ties are broken; the observable region is the feasible set, and the optimal station is its optimizer. The three constraints are heterogeneous and the optimizer may be a one-dimensional continuum, so we use an analytics-first route with numerics only as evidence.

\subsection{Analysis of Problem 2}

Part~2(1) requires a two-station accuracy model, an accuracy index, station selection for T2 under the 30~s limit, and a degeneracy discussion; Part~2(2) minimizes the latest start with three alerts, six slots and five vehicles and outputs the timeline against a nearest-assignment baseline. The two parts share geometry but not objectives, and 2(1) is a conditional accuracy analysis of two-station bearing intersection.

\subsection{Analysis of Problem 3}

Problem~3 requires a propagation model for $\Delta\psi$, $s_{\max}$, the $(a,b)$ constraints at 80~m, justification of the approximation, comparison of the three error sources, and deployment implications. The difficulties are the different dimensions of the sources, the nonlinear $\operatorname{atan2}$, the two-sided criterion, and the monotonicity needed for the $(a,b)$ constraints.

\subsection{Analysis of Problem 4}

Problem~4 takes the time-optimal 10.355~s plan as the baseline and considers, at $t=5$~s, new alerts, one vehicle failure, communication outages and their combinations, evaluated by $(\eta,t_{\mathrm{last}})$. Disturbances occur mid-execution (locking, outage-triggered rescheduling, continuous new-alert parameters), so the comparison includes a simple-repair baseline to separate reaction from re-optimization.

\subsection{Algorithm Routes and Selection}

Table~\ref{tab:2-2} gives a compact route summary; full details are in Appendix~B.1--B.2.

\begin{table}[htbp]
\centering
\caption{Compact algorithm routes for the four problems}
\label{tab:2-2}
{\small
\renewcommand{\arraystretch}{1.08}
\begin{tabularnb}{@{}p{12mm}p{64mm}p{74mm}@{}}
  \toprule
  Problem & Method & Why \\
  \midrule
  Q1 & Closed-form $r\ge\rmin$ + $\rmin$-circle sampling + 3-D slab/visibility graph + lexicographic tie-break &
  $Q=1/r$ is strictly monotone, so the optimum is certified in closed form; the optimal set is a one-dimensional arc \\
  \midrule
  Q2(1) & Covariance propagation $\Sigma=A^{-1}DA^{-T}$ + QM--AM bound + $C(5,2)$ pair enumeration + polar candidates/Nelder--Mead &
  The equality condition reduces optimality to a $90^\circ$ pairing test on the $\rmin$ circle; only 10 pairs exist \\
  \midrule
  Q2(2) & 1,000 structures / 5,880 reuse orders + fixed-point timeline + coordinate descent/Nelder--Mead; reuse-aware LB &
  5 vehicles and 6 slots permit exhaustive deterministic enumeration; stations are locally refined; certified gap $28.6\%$ \\
  \midrule
  Q3 & Closed-form first-order propagation + two-sided quantile + $256\times256$ Gauss--Hermite quadrature + MC verification &
  $\sigma_\psi(s)$ is monotone, so the $P$ criterion is invertible; quadrature has no sampling error \\
  \midrule
  Q4 & Event-triggered snapshot + reuse of the Q2(2) solver + 74-scenario enumeration with keep/simple-repair/re-optimization comparison &
  The post-disturbance scale matches Q2(2), so the same solver is reused and all scenarios are comparable \\
  \bottomrule
\end{tabularnb}
}
\end{table}
